% reunion_27_feb.tex

\begin{frame}
    \begin{itemize}
        \item \cite{2024_Yi_Zhang}
        \item \cite{Guo_2024_CCDF}
        \item \cite{lavandier_selective_2021}
        \item \cite{mateos_multiobjective_2016}
    \end{itemize}
\end{frame}
\begin{frame}{\parencite{2024_Yi_Zhang} }{\citetitle{2024_Yi_Zhang}}
    \begin{itemize}
        \item La creciente demanda de tráfico aéreo genera congestión en las Áreas de Maniobra Terminal (TMA).
        \item Controladores de tráfico aéreo (ATCOs) dependen de decisiones manuales con soporte limitado de automatización.
        \item Se propone una estrategia de ruteo y programación de aeronaves para minimizar conflictos y optimizar la eficiencia en la zona TMA.
        \item \red{Problema uni-objetivo}
    \end{itemize}
\end{frame}

\begin{frame}{\parencite{2024_Yi_Zhang} }{\citetitle{2024_Yi_Zhang}}
    \textbf{Modelado y Algoritmos:}
    \begin{itemize}
        \item Uso de rutas de llegada estándar (STARs) como restricciones.
        \item Modelo mixto lógico con optimización de separación mínima y orden de aterrizaje.
        \item Métodos de solución:
              \begin{itemize}
                  \item \textbf{Programación Entera Mixta (MILP):} resolución exacta con Gurobi.
                  \item \textbf{Búsqueda de Caminos Multi-Agente (MAPF):} planificación de rutas sin conflictos con programación por restricciones.
                  \item \bfb{Metaheurística (HS):} Algoritmo Harmony Search para mejorar eficiencia computacional.
              \end{itemize}
    \end{itemize}
\end{frame}

\begin{frame}{\parencite{2024_Yi_Zhang} }{\citetitle{2024_Yi_Zhang}}
    \bfb{Función Objetivo}
    \begin{itemize}
        \item \textbf{Minimización del tiempo de aterrizaje del último avión (Makespan).}
        \item Maximiza el uso de la infraestructura y reduce los tiempos de espera en el aire.
              \[
                  J = \min \left\{\max_{f \in F,\, j \in D} \left\{AR^f_j\right\}\right\}
              \]
              donde \(AR^f_j\) es el tiempo de llegada de la aeronave \(f\) al punto final \(j\).
    \end{itemize}
\end{frame}

\begin{frame}{\parencite{2024_Yi_Zhang} }{\citetitle{2024_Yi_Zhang}}
    \textbf{Restricciones principales:}
    \begin{itemize}
        \item \textbf{Asignación de Entrada:} Cada aeronave se asigna a un punto de entrada en TMA.
        \item \textbf{Conservación de Flujo:} Cada aeronave sigue una trayectoria continua sin desvíos innecesarios.
        \item \textbf{Limitaciones de Conexión y Espera:}
              \begin{itemize}
                  \item Solo se permiten esperas en puntos específicos.
                  \item Selección de velocidad en cada tramo de la ruta.
              \end{itemize}
        \item \textbf{Evitar Conflictos:} Separación mínima de tiempo entre aeronaves en cada punto de la trayectoria.
        \item \textbf{Prohibición de Adelantamientos:} Una aeronave no puede superar a otra en la misma ruta.
    \end{itemize}
\end{frame}



\begin{frame}{\parencite{2024_Yi_Zhang} }{\citetitle{2024_Yi_Zhang}}
    \bfb{Resultados y Simulaciones}
    \begin{itemize}
        \item Se probaron 15 escenarios de diferentes escalas y datos reales de ADS-B en Singapur.
        \item Reducción del tiempo total de aterrizaje en hasta \bfb{10 minutos} por aeronave.
        \item Reducción de más de \bfb{80 minutos} en el tiempo de vuelo total.
        \item MILP y MAPF tienen mejor desempeño en problemas pequeños, mientras que \bfb{HS es más escalable}.
    \end{itemize}
\end{frame}
\begin{frame}{\parencite{2024_Yi_Zhang} }{\citetitle{2024_Yi_Zhang}}
    \begin{figure}
        \centering
        \includegraphics[width=0.95\textwidth]{2024_zhang_A_Study}
        \caption*{\tiny\citeauthor{2024_Yi_Zhang}; \citetitle{2024_Yi_Zhang}, \citeyear{2024_Yi_Zhang}}
    \end{figure}
\end{frame}

\begin{frame}{\parencite{2024_Yi_Zhang} }{\citetitle{2024_Yi_Zhang}}
    \bfb{Características de la Solución HS}
    \begin{itemize}
        \item Cada solución se compone de:
              \begin{itemize}
                  \item \textbf{Ruta (\(X^r_f\))}: Secuencia de waypoints (por ejemplo, 2--5--8--11--15--Runway).
                  \item \textbf{Holding (\(X^h_f\))}: Tiempo de espera asignado en cada waypoint donde se permite holding.
                  \item \textbf{Velocidad (\(X^s_f\))}: Selección de nivel de velocidad (0,1,2,3) para cada tramo.
              \end{itemize}
        \item La solución se genera de forma aleatoria pero respetando las restricciones duras del problema.
        \item La evaluación se realiza combinando el costo (makespan) y penalizaciones por violaciones a restricciones.
    \end{itemize}
\end{frame}

\begin{frame}{\parencite{2024_Yi_Zhang} }{\citetitle{2024_Yi_Zhang}}
    \bfb{Ejemplo de Solución generada aleatoriamente para Aeronave despachada desde Waypoint 2}
    \vspace{0.5em}
    \begin{figure}
        \centering
        \includegraphics[width=0.9\textwidth]{table_2_zhang}
        \caption*{\parencite{2024_Yi_Zhang}}
    \end{figure}
\end{frame}


\begin{frame}{Conclusiones}
    \begin{itemize}
        \item La estrategia propuesta mejora la eficiencia y seguridad en TMA.
        \item El enfoque híbrido combina exactitud (MILP), rapidez (MAPF) y escalabilidad (HS).
        \item Futuras mejoras pueden incluir factores meteorológicos y múltiples pistas.
    \end{itemize}
\end{frame}

% % ============== ============== ============== ============== ============== ============== ==============

% \section{Paper 2: Cooperative Co-Evolution for Large-Scale Multiobjective ATFM}
\begin{frame}{\parencite{Guo_2024_CCDF} }{\citetitle{Guo_2024_CCDF}}
    \begin{itemize}
        \item \textbf{Objetivo General:}
              Optimizar el flujo de tráfico aéreo minimizando retrasos y conflictos.
        \item \textbf{Funciones Objetivo:}
              \begin{itemize}
                  \item \(\min TFD\): Minimizar el retraso total de vuelo, definido como la suma de las diferencias entre el tiempo de despegue real y el esperado, y el tiempo de llegada real y el esperado.
                  \item \(\min TNC\): Minimizar el número total de conflictos entre trayectorias, detectados cuando la distancia entre aeronaves es inferior a un umbral (por ejemplo, 5 millas náuticas).
              \end{itemize}
    \end{itemize}
\end{frame}
\begin{frame}{\parencite{Guo_2024_CCDF} }{\citetitle{Guo_2024_CCDF}}
    \begin{itemize}
        \item \textbf{Restricciones Clave:}
              \begin{itemize}
                  \item Cada aeronave debe seleccionar una ruta de entre un conjunto predefinido.
                  \item Las variables de tiempo (despegue) y velocidad deben mantenerse dentro de rangos permitidos.
                  \item Restricciones de capacidad:
                        \begin{itemize}
                            \item Capacidad de despegue y aterrizaje en aeropuertos.
                            \item Capacidad máxima en sectores del espacio aéreo.
                        \end{itemize}
                  \item Garantizar la separación mínima entre aeronaves en el tiempo y el espacio.
              \end{itemize}
    \end{itemize}
\end{frame}

\begin{frame}{\parencite{Guo_2024_CCDF} }{\citetitle{Guo_2024_CCDF}}
    \begin{itemize}
        \item \textbf{Metodología:}
              Se utiliza un enfoque de \textbf{Co-Evolución Cooperativa (CC)} para ATFM a gran escala, que incluye:
              \begin{itemize}
                  \item \textbf{Descomposición Fuzzy:} Divide el problema original en subcomponentes de menor dimensión, aprovechando las correlaciones espaciotemporales entre aeronaves.
                  \item \textbf{Optimización de Subcomponentes:} Cada subproblema se resuelve mediante un metaheurístico (por ejemplo, un algoritmo evolutivo) y se coopera mediante un archivo externo que recoge soluciones no dominadas.
                  \item \textbf{Asignación Probabilística de Recursos:} Se asignan evaluaciones adicionales a los subcomponentes con mayor contribución, medida mediante el incremento del hipervolumen (HV) de la solución.
              \end{itemize}
        \item \textbf{Representación de Individuos:}
              Cada individuo representa un \textbf{plan de vuelo} de la forma:
              \[
                  \langle t_i, r_i, v_i \rangle
              \]
              donde:
              \begin{itemize}
                  \item \(t_i\): Tiempo de despegue de la aeronave \(i\) (ajustable dentro de un intervalo).
                  \item \(r_i\): Ruta seleccionada, definida como una secuencia de waypoints.
                  \item \(v_i\): Perfil de velocidad, en el que cada tramo de la ruta se asigna una velocidad elegida de un conjunto discreto.
              \end{itemize}
    \end{itemize}
\end{frame}

\begin{frame}{\parencite{Guo_2024_CCDF} }{\citetitle{Guo_2024_CCDF}}
    \begin{itemize}
        \item \textbf{Framework Global:}
              \begin{itemize}
                  \item Inicialización de la población y cálculo de los objetivos.
                  \item Descomposición dinámica del problema basada en un \red{knee point} de las soluciones no dominadas, que sirve de guía para la cooperación.
              \end{itemize}
        \item \textbf{Optimización de Subcomponentes:}
              \begin{itemize}
                  \item Cada subpoblación evoluciona de forma colaborativa utilizando un algoritmo evolutivo.
                  \item Las mejoras se miden mediante la contribución al hipervolumen (HV) del conjunto de soluciones.
              \end{itemize}
        \item \textbf{Asignación de Recursos:}
              \begin{itemize}
                  \item Se calcula una probabilidad de selección para cada subcomponente según su contribución.
                  \item Se emplea un mecanismo de ruleta para asignar evaluaciones adicionales, permitiendo a subcomponentes estancados tener nuevas oportunidades.
              \end{itemize}
    \end{itemize}
\end{frame}


% % ============== ============== ============== ============== ============== ============== ==============
%---------------------------------------
% \section{Paper 1: Multiobjective Simulated Annealing for Collision Avoidance in ATM}
\begin{frame}{\parencite{mateos_multiobjective_2016} }{\citetitle{mateos_multiobjective_2016}}
    \begin{itemize}
        \item \textbf{Objetivo General:}
              Minimizar múltiples criterios relacionados con la evitación de colisiones.
        \item \textbf{Funciones Objetivo:}
              \begin{itemize}
                  \item Minimizar la magnitud promedio de las maniobras y su dispersión.
                  \item Minimizar el riesgo de colisión (medido a partir de la proximidad entre cilindros de seguridad).
                  \item Minimizar el número de maniobras, lo que reduce la carga de trabajo de los controladores.
                  \item Minimizar retrasos temporales y desviaciones en los puntos de salida.
              \end{itemize}
        \item \textbf{Restricciones:}
              \begin{itemize}
                  \item \textbf{Seguridad:} Los “cilindros de seguridad” (definidos por distancias mínimas horizontales y verticales) no deben superponerse.
                  \item Límites en velocidades, altitudes y ángulos de giro (para maniobras de cambio de velocidad, giro y altitud).
              \end{itemize}
        \item \textbf{Metodología:}
              Se utiliza un algoritmo de \textbf{Simulated Annealing} multiobjetivo (AMOSA) que explora el espacio de soluciones mediante una estrategia de enfriamiento.
    \end{itemize}
\end{frame}

\begin{frame}{\parencite{mateos_multiobjective_2016} }{\citetitle{mateos_multiobjective_2016}}
    \begin{block}{Representación de Individuos}
        Cada individuo es un vector que, para cada aeronave, contiene:
        \begin{itemize}
            \item \textbf{Decisión de maniobra:} Variables binarias \(vc, ac, tc\) que indican si se realiza un cambio de velocidad, altitud o giro, respectivamente.
            \item \textbf{Magnitud de la maniobra:} Un valor continuo \(\Delta\) en un rango predefinido (por ejemplo, \([-1,1]\)), que representa la intensidad del cambio.
        \end{itemize}
    \end{block}
    \begin{block}{Funciones Objetivo (Ejemplos)}
        \begin{itemize}
            \item \(\min f_1 =\) Promedio de magnitud de maniobras \(+\) dispersión entre ellas.
            \item \(\min f_2 =\) Riesgo de colisión, basado en las intersecciones de los cilindros de seguridad.
            \item \(\min f_3 =\) Número total de maniobras realizadas (con la restricción \(vc+ac+tc \leq 1\) por aeronave).
            \item \(\min f_4 =\) Retrasos en tiempos de salida/entrada.
            \item \(\min f_5 =\) Desviación en el punto de salida respecto al plan original.
        \end{itemize}
    \end{block}
\end{frame}

%---------------------------------------

\begin{frame}{\parencite{mateos_multiobjective_2016} }{\citetitle{mateos_multiobjective_2016}}
    \begin{itemize}
        \item Ambos trabajos enfrentan problemas de optimización en ATM/ATFM con múltiples objetivos y restricciones complejas.
        \item \textbf{\parencite{mateos_multiobjective_2016}}
              \begin{itemize}
                  \item Se centra en la resolución de conflictos mediante maniobras individuales.
                  \item Representa a cada aeronave mediante un vector de decisiones de maniobra.
              \end{itemize}
        \item \textbf{\parencite{Guo_2024_CCDF}:}
              \begin{itemize}
                  \item Se enfoca en la optimización integral del flujo aéreo (horarios, rutas, velocidades).
                  \item Utiliza la descomposición y cooperación entre subcomponentes para manejar problemas a gran escala.
              \end{itemize}
        \item Ambos enfoques destacan el uso de metaheurísticas para explorar grandes espacios de búsqueda en problemas NP-hard, garantizando soluciones cercanas al óptimo en contextos reales.
    \end{itemize}
\end{frame}

% % ============== ============== ============== ============== ============== ============== ==============
\begin{frame}{\parencite{lavandier_selective_2021} }{\citetitle{lavandier_selective_2021}}

\end{frame}




\thispagestyle{plain}
\begin{frame}
    \tiny
    \printbibliography


\end{frame}



